\chapter{Unit 1: Computer Languages}
\label{chap:unit1_qb}

\section{Practice Question Set (50 MCQs)}

\mcqitem{1}{Which of the following is the lowest-level programming language understood directly by a computer's CPU?}{Machine language}{Assembly language}{High-level language}{Natural language}

\mcqitem{2}{What symbols are used to represent data and instructions in machine language?}{Alphabets A to Z}{Binary digits 0 and 1}{Hexadecimal codes only}{ASCII characters}

\mcqitem{3}{A statement in assembly language uses readable abbreviations for operations called:}{Operands}{Mnemonics}{Identifiers}{Macros}

\mcqitem{4}{Which software tool is required to convert assembly language programs into machine code?}{Compiler}{Interpreter}{Assembler}{Linker}

\mcqitem{5}{A translator that converts an entire high-level program into machine code before execution is a/an:}{Interpreter}{Compiler}{Assembler}{Text Editor}

\mcqitem{6}{An interpreter translates and executes high-level language programs:}{All at once prior to execution}{Statement-by-statement during runtime}{Only after linking with libraries}{During compile time}

\mcqitem{7}{In an assembly instruction such as \texttt{ADD R1, R2}, the operation code \texttt{ADD} is known as the:}{Operand}{Opcode}{Mnemonic address}{Label}

\mcqitem{8}{Which of the following programming languages is classified as a high-level language?}{Machine Code}{C++}{Assembly Language}{Microcode}

\mcqitem{9}{In an assembly instruction \texttt{MOV R1, \#5}, what does the symbol \texttt{\#5} represent?}{An opcode}{An immediate operand value}{A memory address pointer}{A condition flag}

\mcqitem{10}{What is the primary advantage of high-level programming languages over machine language?}{They execute without a CPU}{They are portable and human-readable}{They require no translation}{They directly control CPU registers}

\mcqitem{11}{Which phase of compilation groups raw characters into meaningful tokens such as identifiers and keywords?}{Syntax analysis}{Lexical analysis}{Semantic analysis}{Code optimization}

\mcqitem{12}{Checking whether statements in a source code conform to the grammatical rules of the language is done during:}{Lexical analysis}{Syntax analysis}{Code generation}{Linking}

\mcqitem{13}{Detecting a type mismatch error (such as assigning a string to an integer variable) is the responsibility of:}{Semantic analysis}{Lexical analysis}{The loader}{The preprocessor}

\mcqitem{14}{A syntax error in a C++ program is detected at:}{Compile time}{Run time}{Link time}{Load time}

\mcqitem{15}{What is the role of a preprocessor in the C/C++ compilation pipeline?}{To execute machine code}{To expand macros and include header files}{To link object files}{To load the program into RAM}

\mcqitem{16}{The software that combines multiple compiled object files and standard library routines into an executable is the:}{Assembler}{Compiler}{Linker}{Loader}

\mcqitem{17}{Which program is responsible for bringing the executable file from secondary storage into primary memory (RAM) for execution?}{Compiler}{Linker}{Loader}{Interpreter}

\mcqitem{18}{The four stages of the CPU instruction execution cycle in correct order are:}{Decode $\rightarrow$ Fetch $\rightarrow$ Store $\rightarrow$ Execute}{Fetch $\rightarrow$ Decode $\rightarrow$ Execute $\rightarrow$ Store}{Execute $\rightarrow$ Fetch $\rightarrow$ Decode $\rightarrow$ Store}{Fetch $\rightarrow$ Execute $\rightarrow$ Decode $\rightarrow$ Store}

\mcqitem{19}{In an assembler, what is the primary task performed during the first pass of a two-pass assembler?}{Generating final binary machine instructions}{Recording symbol and label definitions with their memory addresses}{Linking object files}{Loading the program into RAM}

\mcqitem{20}{What does a two-pass assembler do during its second pass?}{Translates mnemonics and symbolic references into machine code}{Identifies syntax errors only}{Allocates heap memory}{Runs the code}

\mcqitem{21}{Which type of programming error allows a program to compile and run successfully but produces incorrect output?}{Syntax error}{Logic error}{Linker error}{Lexical error}

\mcqitem{22}{Division by zero occurring during program execution is classified as a:}{Syntax error}{Compile-time error}{Runtime error}{Semantic error}

\mcqitem{23}{The organized sequence used to transform a real-world problem into a working software application is called the:}{Instruction cycle}{Program development cycle}{Compilation pipeline}{Bus arbitration cycle}

\mcqitem{24}{What is the first step in the program development cycle?}{Writing program code}{Problem definition}{Program maintenance}{Compiling the code}

\mcqitem{25}{A step-by-step graphical representation of an algorithm using standard geometric symbols is called a:}{Pseudocode}{Flowchart}{Syntax tree}{Program draft}

\mcqitem{26}{A finite, ordered set of unambiguous instructions designed to solve a specific problem is known as a/an:}{Operand}{Algorithm}{Token}{Mnemonic}

\mcqitem{27}{Testing program execution using edge values such as \texttt{hours = 0} is an example of:}{Syntax verification}{Boundary value testing}{Preprocessing}{Lexical scanning}

\mcqitem{28}{Why is machine language considered hardware-dependent?}{Because it uses English keywords}{Because each CPU family has its own unique binary instruction set}{Because it requires a high-level interpreter}{Because it operates only on laptops}

\mcqitem{29}{Which language requires an assembler for conversion to machine code?}{Java}{Assembly language}{Python}{C}

\mcqitem{30}{Which tool provides an interactive environment with line-by-line execution and immediate feedback, ideal for scripting and beginners?}{Compiler}{Interpreter}{Linker}{Assembler}

\mcqitem{31}{What happens during the \textit{Fetch} stage of the CPU instruction cycle?}{The ALU performs an addition}{The instruction is read from memory into a CPU register}{The preprocessor expands macros}{The linker connects libraries}

\mcqitem{32}{What happens during the \textit{Decode} stage of the CPU instruction cycle?}{The CPU determines the operation to be performed and the operands involved}{The result is saved to a hard drive}{The program is recompiled}{The operating system restarts}

\mcqitem{33}{In the statement \texttt{total = price * quantity}, what provides the abstraction that hides CPU registers from the programmer?}{The high-level language translator}{The physical power supply}{The motherboard bus}{The keyboard controller}

\mcqitem{34}{Which phase of the program development cycle involves modifying software after release to correct newly discovered bugs or adapt to new hardware?}{Problem Analysis}{Maintenance}{Coding}{Algorithm Design}

\mcqitem{35}{What does source code mean?}{Binary instructions executed by the CPU}{Human-readable program text written by a programmer}{Machine language saved on a disk}{Electrical signals on a motherboard}

\mcqitem{36}{Which of the following file extensions typically contains high-level C programming source code?}{\texttt{.c}}{\texttt{.obj}}{\texttt{.exe}}{\texttt{.bin}}

\mcqitem{37}{What is an object file?}{A high-level source text file}{A compiled file containing machine instructions and symbol tables prior to linking}{An image file created by an image editor}{A document explaining software requirements}

\mcqitem{38}{Why do compiled programs generally execute faster than interpreted programs?}{They do not use RAM}{Translation to machine code has already been completed prior to execution}{They avoid using CPU registers}{They have no runtime errors}

\mcqitem{39}{In assembly language, what is a label used for?}{To specify a high-level data type}{To name a specific memory location or jump target in the code}{To run preprocessor commands}{To exit the operating system}

\mcqitem{40}{Which language typically uses a hybrid approach of compiling to bytecode and then executing on a Virtual Machine?}{Assembly language}{Java}{Pure Machine code}{Fortran I}

\mcqitem{41}{A statement such as \texttt{\#include <stdio.h>} in C is directed to the:}{Preprocessor}{Linker}{Loader}{Operating system kernel}

\mcqitem{42}{The set of formal rules governing the structure and character sequences of valid statements in a language is called:}{Semantics}{Syntax}{Heuristics}{Mnemonics}

\mcqitem{43}{The meaning or logical effect of a syntactically valid programming statement is referred to as its:}{Grammar}{Semantics}{Tokenization}{Format}

\mcqitem{44}{Which of the following is an example of an interpreted scripting language?}{Python}{C}{C++}{Assembly}

\mcqitem{45}{In the instruction \texttt{grossPay = hours * rate}, what are \texttt{hours} and \texttt{rate}?}{Output variables}{Input operands}{Opcodes}{Directives}

\mcqitem{46}{What type of error is caused by misspelling a language keyword such as writing \texttt{whille} instead of \texttt{while}?}{Runtime error}{Syntax error}{Logic error}{Linker error}

\mcqitem{47}{When a compiler replaces repetitive code with faster equivalent machine instructions without changing program output, it is performing:}{Lexical analysis}{Code optimization}{Semantic checking}{Linking}

\mcqitem{48}{A program that translates symbolic assembly language instructions into machine language instructions on a one-to-one basis is an:}{Interpreter}{Assembler}{Operating system}{Compiler}

\mcqitem{49}{What is pseudocode?}{An executable binary script}{An informal, English-like description of an algorithm's steps}{A hardware test program}{An assembly directive}

\mcqitem{50}{Which component physically performs the arithmetic addition when the CPU executes an addition instruction?}{Control Unit}{Arithmetic Logic Unit (ALU)}{Program Counter}{Secondary Storage}

\newpage
\section{Answer Key \& Explanations --- Unit 1}

\begin{answerkeytable}{Unit 1: Computer Languages --- Answer Key}
\begin{tabular}{*{10}{c}}
\toprule
\textbf{Q} & \textbf{Ans} & \textbf{Q} & \textbf{Ans} & \textbf{Q} & \textbf{Ans} & \textbf{Q} & \textbf{Ans} & \textbf{Q} & \textbf{Ans} \\
\midrule
1 & \textbf{A} & 11 & \textbf{B} & 21 & \textbf{B} & 31 & \textbf{B} & 41 & \textbf{A} \\
2 & \textbf{B} & 12 & \textbf{B} & 22 & \textbf{C} & 32 & \textbf{A} & 42 & \textbf{B} \\
3 & \textbf{B} & 13 & \textbf{A} & 23 & \textbf{B} & 33 & \textbf{A} & 43 & \textbf{B} \\
4 & \textbf{C} & 14 & \textbf{A} & 24 & \textbf{B} & 34 & \textbf{B} & 44 & \textbf{A} \\
5 & \textbf{B} & 15 & \textbf{B} & 25 & \textbf{B} & 35 & \textbf{B} & 45 & \textbf{B} \\
6 & \textbf{B} & 16 & \textbf{C} & 26 & \textbf{B} & 36 & \textbf{A} & 46 & \textbf{B} \\
7 & \textbf{B} & 17 & \textbf{C} & 27 & \textbf{B} & 37 & \textbf{B} & 47 & \textbf{B} \\
8 & \textbf{B} & 18 & \textbf{B} & 28 & \textbf{B} & 38 & \textbf{B} & 48 & \textbf{B} \\
9 & \textbf{B} & 19 & \textbf{B} & 29 & \textbf{B} & 39 & \textbf{B} & 49 & \textbf{B} \\
10 & \textbf{B} & 20 & \textbf{A} & 30 & \textbf{B} & 40 & \textbf{B} & 50 & \textbf{B} \\
\bottomrule
\end{tabular}
\end{answerkeytable}

\begin{expbox}[Selected Question Explanations]
\begin{itemize}[leftmargin=5mm]
    \item \textbf{Q.1 (A):} Machine language is written purely in binary and is the native language directly understood and processed by the CPU.
    \item \textbf{Q.3 (B):} Assembly language replaces tedious binary bit sequences with readable letter codes called mnemonics (e.g., \texttt{MOV}, \texttt{ADD}).
    \item \textbf{Q.5 (B):} A compiler translates the entire source code into machine or object code in a separate compilation step prior to execution.
    \item \textbf{Q.6 (B):} An interpreter processes code line-by-line during runtime, stopping immediately when an error is encountered.
    \item \textbf{Q.11 (B):} Lexical analysis (the scanner) reads raw character streams and groups them into lexical tokens like identifiers, literals, and keywords.
    \item \textbf{Q.16 (C):} The linker connects independent object files and standard library routines together to produce a single executable binary.
    \item \textbf{Q.18 (B):} The CPU instruction cycle always runs in the fundamental order: Fetch $\rightarrow$ Decode $\rightarrow$ Execute $\rightarrow$ Store.
    \item \textbf{Q.21 (B):} Logic errors do not violate grammar rules so the program compiles and executes, but it produces incorrect mathematical or logical results.
    \item \textbf{Q.38 (B):} Compiled binaries execute faster because all translation overhead occurs beforehand, unlike interpreted code which requires runtime evaluation.
\end{itemize}
\end{expbox}

\newpage
